% 176_T0_Shor_Photon_De_ch.tex -- Kapitelinhalt Dok. 176
% Aufbereitung und Transformation: zwei getrennte Vorgänge
% Vollständige Neufassung August 2026 (R75)
% Mit Quellenangaben erweitert und Inhalte angepasst

\section{Aufbereitung und Transformation: zwei getrennte Vorgänge}

\subsection*{Zusammenfassung}

Shors Algorithmus zerfällt in zwei Vorgänge, die regelmäßig vermengt
werden, obwohl sie nichts miteinander zu tun haben:

\begin{center}
\begin{tabular}{lll}
\toprule
 & \textbf{Aufbereitung} & \textbf{Transformation} \\
\midrule
Eingang & Zahlen $a$, $N$ & aufbereitetes Signal \\
Ausgang & Signal mit $f(x) = a^x \bmod N$ & Frequenzdarstellung \\
Natur & klassische Arithmetik & Wellenoperation \\
Aufwand & \textbf{teuer} ($O(n^3)$ Gatter) & \textbf{billig} ($O(n^2)$) \\
\bottomrule
\end{tabular}
\end{center}

Die Aufbereitung ist schwierig, weil die modulare Exponentiation
schwierig ist -- in jeder Technologie gleichermaßen. Sie hängt
\emph{nicht} an quantenmechanischen Eigenschaften. Die Transformation
ist in jeder Technologie dieselbe Wellenoperation und in jeder billig.
Der Unterschied zwischen Optik und Qubit-Register liegt weder im einen
noch im anderen, sondern in Betriebstemperatur und
Kopplungsmechanismus.


\begin{keyresult}[Drei getrennte Parameterräume]
Zur Vermeidung von Verwechslungen unterscheidet dieses Dokument strikt:
\begin{enumerate}
  \item \textbf{Geometrischer Raum:} $\xi = 4/30000 \approx 1{,}333\times10^{-4}$
    -- der Strukturparameter der FFGFT. Fest, nicht einstellbar.
  \item \textbf{Algebraischer Raum:} $\xi_{\text{Higgs}} \approx
    1{,}038\times10^{-5}$ -- Feldkorrektur im Vakuumerwartungswert.
  \item \textbf{Zahlentheoretischer Raum} $\mathbb{Z}/N\mathbb{Z}$:
    \emph{kein} universeller Parameter. Die Periode $r$ wird durch
    $\mathrm{lcm}(p-1,q-1)$ bestimmt -- eine individuelle arithmetische
    Eigenschaft von $N$.
\end{enumerate}
Die Faktorisierung lebt ausschließlich im dritten Raum. Weder $\xi$ noch
$\xi_{\text{Higgs}}$ haben dort eine Funktion.
\end{keyresult}

\section{Die Aufbereitung}
\label{sec:aufbereitung}

\textbf{Eingang:} die Zahlen $a$ und $N$.
\textbf{Ausgang:} ein Zustand oder Signal, das die Folge
$f(x) = a^x \bmod N$ trägt.

\subsection{Was hier geschieht}

\textbf{[K]} Die modulare Exponentiation ist eine \emph{klassische,
deterministische} Funktion. Ihre Auswertung enthält keinen
quantenmechanischen Anteil: Auf einem Quantenrechner wird sie als
reversibler klassischer Schaltkreis ausgeführt \cite{Vedral1996}, in einem
optischen Aufbau als räumliche oder spektrale Struktur eingeprägt.

\subsection{Die Kosten}

\textbf{[K]} Der Aufwand hängt an der Arithmetik, nicht an der
gewählten Technologie:

\begin{center}
\begin{tabular}{lll}
\toprule
Realisierung & Aufwand & für RSA-2048 \\
\midrule
Quantenschaltkreis & $O(n^3)$ Gatter & $8{,}6\cdot10^{9}$ Gatter \\
Optische Maske & $2^n$ Elemente & $10^{617}$ Elemente \\
Klassische Tabelle & $2^n$ Werte & $10^{617}$ Werte \\
\bottomrule
\end{tabular}
\end{center}

Der Quantenschaltkreis ist die günstigste bekannte Variante: Er
kodiert die Rechenvorschrift statt der Werte und kommt daher mit
polynomiellem Aufwand aus. Eine optische Maske müsste dagegen jeden
Funktionswert einzeln tragen -- für RSA-2048 mehr Elemente, als das
beobachtbare Universum Atome hat ($10^{80}$).

\begin{keyresult}[Die Aufbereitung ist ein arithmetisches Problem]
Sie ist teuer, weil die modulare Exponentiation teuer ist -- und das
gilt in jeder Technologie. Ein Verfahren, das sie zu umgehen
verspricht, verschiebt den Aufwand, statt ihn zu beseitigen.
\textbf{Die Aufbereitung hängt nicht an quantenmechanischen
Eigenschaften.}
\end{keyresult}

\subsection{Was die Superposition dabei leistet}

\textbf{[X]} Nichts. Der Zustand
$\sum_x \ket{x}\ket{a^x \bmod N}$ enthält zwar alle Paare, aber eine
Messung liefert genau eines -- zufällig gezogen. Danach weiß man so
viel wie nach einer einzigen klassischen Auswertung. Die Superposition
vereinfacht die Arithmetik nicht; sie legt denselben Schaltkreis über
alle Eingaben.

Barenco et al. \cite{Barenco1995} zeigen, dass Quantengatter elementare
arithmetische Operationen (Addition, Multiplikation) ausführen, die in
ihrer Komplexität klassischen Operationen entsprechen. Die Superposition
ermöglicht paralleles Anlegen, aber nicht paralleles Verarbeiten der
Arithmetik -- jeder Rechenschritt muss für jede Komponente der
Superposition einzeln durchgeführt werden.

\subsection{Die modulare Exponentiation als Flaschenhals}

Die \emph{modulare Exponentiation} wird in der Forschung einhellig als der
Flaschenhals (das ``bottleneck'') von Shors Algorithmus bezeichnet
\cite{Vedral1996}. Vedral, Barenco und Ekert analysieren die Komplexität
von Quanten-Addierern und Multiplizierern und dokumentieren den hohen
Aufwand für die modulare Exponentiation.

\begin{center}
\begin{tabular}{llll}
\toprule
Architektur & Qubits & Quantenkosten & Quelle \\
\midrule
Pavlidis \& Gizopoulos & $\sim 9n + 2$ & $\sim 1600 n^3$ Gatter & \cite{Pavlidis2014} \\
Ragavan & $\sim 10.4n$ & $\sim 45.7 n^{1/2}$ Multiplikationen & \cite{Ragavan2018} \\
Patoary et al. (DFT-basiert) & -- & $O(n^3)$ Gatter & \cite{Patoary2021} \\
\bottomrule
\end{tabular}
\end{center}

Die Kosten sind polynomiell ($n^3$) oder sogar besser -- aber sie
existieren und sind der dominierende Faktor \cite{Nielsen2010}.

Ein aktueller Forschungsansatz von Ekerå und Gidney \cite{Ekera2023}
beweist sogar, dass man einen Teil dieser Quanten-Arithmetik durch
\emph{klassische Vorab-Berechnung} ersetzen kann. Die Idee: Man berechnet
eine Tabelle von Werten klassisch und speichert sie, um sie dann im
Quanten-Schaltkreis abzurufen. Das reduziert die Anzahl der teuren
Quanten-Operationen um einen Faktor von bis zu $\omega$. Das ist der
stärkste Beweis dafür, dass die Qubits die Arithmetik nicht
beschleunigen -- man kann sie sogar durch klassische Vorarbeit ersetzen,
um die Quantenschaltung zu entlasten.

\section{Die Fouriertransformation}
\label{sec:ft}

\textbf{Eingang:} das aufbereitete Signal.
\textbf{Ausgang:} dessen Frequenzdarstellung.

\subsection{Was hier geschieht}

\textbf{[K]} Die Fouriertransformation ist eine \emph{Wellenoperation}.
Sie interessiert sich nicht dafür, woher das Signal stammt oder was
es bedeutet -- sie übersetzt eine Darstellung in eine andere. Neue
Information entsteht dabei nicht: Was im Eingangssignal fehlt,
erscheint auch in der Ausgabe nicht.

\subsection{Die Kosten}

\textbf{[K]} Sie ist in jeder Realisierung billig:

\begin{center}
\begin{tabular}{lll}
\toprule
Realisierung & Aufwand & Bemerkung \\
\midrule
Quanten-Fouriertransformation & $O(n^2)$ Gatter & sequenziell \\
Optisches MZI-Netz & Durchlaufzeit & simultan, Lichtgeschwindigkeit \\
Klassische FFT & $O(2^n \cdot n)$ & auf explizite Werte \\
\bottomrule
\end{tabular}
\end{center}

Im Quantenregister und im optischen Netz ist der Aufwand gering. Der
Unterschied zwischen beiden liegt in der Betriebstemperatur und im
Kopplungsmechanismus (Abschnitt~\ref{sec:photonik}), nicht in der
Transformation selbst \cite{Nielsen2010}.

\subsection{Wozu sie dient}

Die gesuchte Periode $r$ ist eine \emph{globale} Eigenschaft der Folge:
Kein einzelner Funktionswert trägt sie, man erkennt sie erst im
Vergleich mehrerer Werte. Die Fouriertransformation macht diesen
Vergleich physikalisch -- durch Interferenz löscht sich alles aus
außer den Frequenzen $k/r$, und aus wenigen Messungen folgt $r$ per
Kettenbruchentwicklung.

\begin{keyresult}[Zwei getrennte Vorgänge]
\emph{Aufbereitung:} arithmetisch, teuer, technologieunabhängig
schwierig.
\emph{Transformation:} wellenmechanisch, billig, in jeder Technologie
dieselbe Operation.

Die beiden hängen nicht zusammen. Man kann die Aufbereitung
verbessern, ohne die Transformation zu berühren, und umgekehrt.
\end{keyresult}

\section{Photonische Realisierung}
\label{sec:photonik}

\subsection{MZI-Netzwerke}

Ein Mach-Zehnder-Interferometer realisiert Phasendrehungen
\begin{equation}
  \Delta\phi = \frac{2\pi jk}{n},
\end{equation}
die exakt den Twiddle-Faktoren der diskreten Fouriertransformation
entsprechen. Ein Netzwerk solcher Interferometer führt die
Fouriertransformation als direkte Wellenoperation aus.

\textbf{[B]} Damit ist die Fouriertransformation selbst kein
Alleinstellungsmerkmal des Quantenrechners. Sie ist eine Wellenoperation, die
sich optisch, akustisch oder elektronisch realisieren lässt. Der Unterschied
liegt in der Zustandspräparation, nicht in der Transformation \cite{Nielsen2010}.

\subsection{Der thermische Vorteil optischer Realisierung}

\textbf{[K]} Ob ein Mode ohne Kühlung im Grundzustand ist, entscheidet das
Verhältnis $h\nu/(k_B T)$. Die mittlere thermische Besetzung beträgt
\begin{equation}
  \bar n = \frac{1}{e^{h\nu/k_B T} - 1}.
\end{equation}

\begin{center}
\begin{tabular}{lrrrr}
\toprule
System & $\nu$ (Hz) & $T$ (K) & $h\nu/k_BT$ & $\bar n$ \\
\midrule
Supraleitendes Qubit & $5{,}0\cdot10^{9}$ & $300$ & $8{,}0\cdot10^{-4}$
  & $1{,}25\cdot10^{3}$ \\
Supraleitendes Qubit & $5{,}0\cdot10^{9}$ & $0{,}020$ & $12{,}0$
  & $6{,}2\cdot10^{-6}$ \\
Photon 1550\,nm & $1{,}93\cdot10^{14}$ & $300$ & $30{,}9$
  & $3{,}7\cdot10^{-14}$ \\
Photon 600\,nm & $5{,}0\cdot10^{14}$ & $300$ & $80{,}0$
  & $1{,}8\cdot10^{-35}$ \\
\bottomrule
\end{tabular}
\end{center}

Bei $5$~GHz und Raumtemperatur sitzen rund \textbf{1250 thermische
Anregungen} im Mode -- das System ist vollständig durchmischt. Für eine
Besetzung unter $0{,}01$ ist $T < 52$~mK erforderlich. Bei optischen
Frequenzen ist der Grundzustand bei Raumtemperatur praktisch exakt besetzt.

\textbf{Kühlkosten.} Die Carnot-Grenze verlangt mindestens $3\cdot10^{4}$
Arbeitseinheiten je bei $10$~mK abgeführter Wärmeeinheit; real deutlich mehr.
Ein Verdünnungskryostat zieht $10$--$25$~kW im Dauerbetrieb, ein optisches
Netz nichts.

\begin{keyresult}[Thermodynamisch privilegiert, nicht rechnerisch]
Der thermische Vorteil folgt direkt aus $h\nu/k_BT$ und ist der stärkste
bekannte Vorteil photonischer gegenüber supraleitender Realisierung. Er
betrifft \emph{Infrastruktur} -- Kühlung, Energie, Baugröße --, nicht die
Rechenmächtigkeit. Beide Realisierungen können prinzipiell dieselben
Algorithmen ausführen; die Photonik hat eigene Schwierigkeiten
(deterministische Zweiphotonen-Gatter, Photonenverluste als dominanter
Fehlerkanal).
\end{keyresult}

\subsection{Woher die exponentielle Dimension kommt}

\textbf{[K]} Eine Gegenüberstellung ``Optik gegen Quantenrechner''
wäre unzutreffend -- sie stellte eine Technologie gegen ein Rechenmodell.
Photonische Quantenrechner existieren (KLM-Schema, messungsbasierte
Verfahren). Entscheidend ist allein die Struktur des Zustandsraums:

\begin{center}
\begin{tabular}{lr}
\toprule
Aufbau & Dimension \\
\midrule
klassische Welle in $n$ Moden & $n$ \\
\emph{ein} Photon in $n$ Moden & $n$ \quad(dasselbe) \\
$n$ verschränkte Zweiniveausysteme & $2^n$ \\
$m$ ununterscheidbare Photonen in $n$ Moden & $\binom{n+m-1}{m}$ \\
\bottomrule
\end{tabular}
\end{center}

Ein einzelnes Photon in einem MZI-Netz realisiert dieselbe
$n$-dimensionale Unitäre wie eine klassische Welle -- klassisch
simulierbar. Mit mehreren ununterscheidbaren Photonen wächst die Dimension
exponentiell: 25 Photonen in 50 Moden ergeben $3{,}5\cdot10^{19}$. Der
Sprung kommt aus der \emph{Tensorproduktstruktur}, nicht aus dem
Trägermedium.

\subsection{Was den Rechenvorteil trägt -- offene Frage}

\textbf{[K]} Naheliegend wäre, den Vorteil der Verschränkung zuzuschreiben.
Der Satz von Gottesman-Knill (1998) widerlegt das: Jeder Schaltkreis aus
ausschließlich Clifford-Gattern ($H$, $S$, CNOT) ist klassisch in
Polynomzeit simulierbar, unabhängig von der erzeugten Verschränkung. Der
GHZ-Zustand $(\ket{0\cdots0} + \ket{1\cdots1})/\sqrt2$ ist maximal
verschränkt und entsteht aus $H$ und $n-1$ CNOTs -- bei $1000$ Qubits
genügen rund $2\cdot10^{6}$ Bits für die vollständige klassische
Beschreibung.

Weitere Gegenbelege: DQC1 zeigt vermuteten Vorteil bei nahezu fehlender
Verschränkung; Boson-Sampling beruht auf Ununterscheidbarkeit und
Permanenten. Jozsa und Linden \cite{Jozsa2003} zeigen Verschränkung als
\emph{notwendig} für exponentiellen Vorteil bei reinen Zuständen -- über
Hinreichen sagt der Satz nichts.

\textbf{[X]} Welche Größe den Rechenvorteil trägt, ist nicht geklärt.
Diskutiert werden Kontextualität, Nicht-Stabilisiertheit
(``Magie'') und die Interferenzstruktur selbst; keiner dieser
Kandidaten ist als alleiniger Träger etabliert. Dieses Dokument benennt
die Frage als offen, statt sie zu beantworten.

\textbf{Was daraus für die Faktorisierung folgt.} Die optische
Realisierung löst das Problem nicht: Der limitierende Schritt ist die
Zustandspräparation -- die Auswertung von $a^x \bmod N$ über alle $x$ --,
und die ist arithmetisch. Ein Netz kann eine präparierte Superposition
transformieren; erzeugen kann es sie nicht.

\subsection{Topologische Einschränkung}

\textbf{[K]} Damit eine Welle im Phasenraum eine geschlossene Resonanz bildet,
muss die \emph{logarithmische Windungszahl} $W = \log_2(r)$ rational sein --
nicht das Frequenzverhältnis $r$ selbst:

\begin{center}
\begin{tabular}{llll}
\toprule
Fall & Verhältnis $r$ & Windungszahl $W$ & Bahn \\
\midrule
Reine Quinte & $3/2$ (rational) & $0{,}584963$ (irrational) & offen \\
Temperierte Quinte & $2^{7/12}$ (irrational) & $7/12$ (rational) & geschlossen \\
\bottomrule
\end{tabular}
\end{center}

Ungedämpfte Schwingungen auf reinen Zahlenverhältnissen erzeugen irrationale
Windungszahlen und laufen als offene Spirale auseinander. Ein optisches
Netzwerk benötigt daher ein diskretes, rationalisiertes Phasenraster, um
scharfe Resonanzpunkte zu erzielen -- die Schließung ist ein Artefakt der
Diskretisierung \cite{Dok316}.

\section{Strukturelle Grenzen}

\subsection{Die Weyl-Obstruktion}

\textbf{[K]} Ein Quantenrechner mit $n$ Qubits bildet einen kompakten,
endlichen Zustandsraum ab. Nach dem Weyl-Gesetz wächst dessen Eigenwertdichte
nach einem Potenzgesetz:
\begin{equation}
  N(\lambda) \sim C_d \cdot \lambda^{d/2}.
\end{equation}
Arithmetische Zähldichten wachsen dagegen wie
\begin{equation}
  N(T) \sim T\ln T
\end{equation}
(Riemann-von Mangoldt). $T\ln T$ ist keine Potenz \cite{Dok316}.

\textbf{Folge:} Ein endlicher Zustandsraum kann ein solches Spektrum
prinzipiell nicht vollständig tragen -- er trägt höchstens einen
Anfangsabschnitt. Für unendliche Spektren dieser Dichte wäre eine
\emph{ausgerollte Dilatationsdomäne} erforderlich (Dok.~315, 316), nicht ein
kompakter Resonanzraum.

\begin{keyresult}[Wichtige Abgrenzung]
Diese Obstruktion trifft \textbf{nicht} Shors Algorithmus. Shor sucht
\emph{eine} Periode \emph{einer konkreten} Zahl $N$ -- eine endliche Aufgabe,
für die ein Register mit $2n+3$ Qubits ausreicht. Die Weyl-Obstruktion
betrifft die vollständige Abbildung \emph{unendlicher} arithmetischer
Spektren. Wer sie gegen Shor anführt, überdehnt sie.

Was Shor bei großen $N$ tatsächlich limitiert, ist die Gatterzahl der
modularen Exponentiation (wächst wie $n^3$) und der
Fehlerkorrektur-Overhead (rund $10^3$ physikalische je logischem Qubit).
Für RSA-2048: etwa $4099$ logische Qubits, $8{,}6\times10^9$ Gatter,
$4{,}1\times10^6$ physikalische Qubits \cite{Gidney2021}. Das ist eine
\emph{Ressourcengrenze} -- prinzipiell durch bessere Hardware überwindbar,
anders als eine strukturelle.
\end{keyresult}

\subsection{Primzahlen als Generatoren}

\textbf{[K]} Die tragende Resonanzstruktur beruht auf Primzahlen. Der
Bikohärenztest an 2469 Riemann-Nullstellen (Dok.~316) zeigt Phasenkopplung
genau bei Primzahlpotenzen ($b^2 = 0{,}76$--$0{,}85$) und keine bei
zusammengesetzten Zahlen ($b^2 = 0{,}018$--$0{,}025$) \cite{Dok316}.

Das physikalische Gegenstück ist die Naturtonreihe: In $n f_0$ führen genau
die Primzahl-Ordnungen neue Klangqualitäten ein; $4 = 2\cdot2$, $6 = 2\cdot3$,
$9 = 3\cdot3$ sind Kombinationen bereits vorhandener Achsen.

\section{Einordnung des Quantenrechners}

\textbf{[K]} Shors Algorithmus faktorisiert auf Quantenhardware theoretisch in
$O((\log N)^3)$ \cite{Nielsen2010} -- eine reale Leistung für ein spezifisches
Problem. Daraus folgt jedoch nicht:

\begin{itemize}
  \item \textbf{[X]} \emph{Ein allgemeiner exponentieller Vorteil von
    Quantenrechnern.} Er ist für die überwiegende Mehrheit der Probleme nicht
    nachgewiesen. Shor ist ein Einzelfall, kein Beleg für einen generellen
    Vorteil.
  \item \textbf{[X]} \emph{Ein Vorteil bei der Fouriertransformation selbst.}
    Diese ist eine Wellenoperation und optisch realisierbar. Der Engpass liegt
    in der Zustandspräparation \cite{Nielsen2010}.
  \item \textbf{[X]} \emph{Ein Beitrag zu strukturellen Fragen der
    Zahlentheorie.} Die Weyl-Obstruktion gilt für endliche Zustandsräume
    uneingeschränkt \cite{Dok316}.
\end{itemize}

\textbf{Der Kern:} Das Auslesen einer Periode ist ein spektroskopisches
Messproblem auf präparierten Wellenfeldern. Die Messung findet, was präpariert
wurde; die Präparation ist der teure Teil.

\section{Bilanz}

\begin{center}
\begin{tabular}{ll}
\toprule
Aussage & Status \\
\midrule
QFT ist eine Interferenzoperation & \textbf{[K]} \\
QFT erzeugt keine neue Information & \textbf{[K]} gemessen \\
Windungszahl entscheidet über Schließung & \textbf{[K]} \\
Primzahlen tragen die Resonanzstruktur & \textbf{[K]} bikohärenzgemessen \\
Fouriertransformation optisch realisierbar & \textbf{[B]} \\
Optik thermodynamisch privilegiert & \textbf{[K]} $h\nu/k_BT$ \\
Exponentielle Dimension aus dem Medium & \textbf{[X]} Tensorstruktur \\
Verschränkung trägt den Rechenvorteil & \textbf{[X]} Gottesman-Knill \\
Aufbereitung ist arithmetisch, nicht quantenmechanisch & \textbf{[K]} \\
Aufbereitung und Transformation sind unabhängig & \textbf{[K]} \\
Allgemeiner Quantenvorteil & \textbf{[X]} nicht nachgewiesen \\
Endlicher Raum trägt unendl.\ $T\ln T$-Spektrum & \textbf{[X]} Weyl \\
Weyl-Obstruktion limitiert Shor & \textbf{[X]} Ressourcengrenze \\
$\xi$ hat Funktion in $\mathbb{Z}/N\mathbb{Z}$ & \textbf{[X]} getrennte Räume \\
\bottomrule
\end{tabular}
\end{center}

\textbf{Der tragende Satz.} Die Fouriertransformation ist eine Projektion --
sie findet Struktur, erzeugt sie aber nicht. Das gilt gleichermaßen für die
klassische FFT, die optische Realisierung und die QFT. Die Faktorisierung lebt
im Raum $\mathbb{Z}/N\mathbb{Z}$, in dem die Geometrieparameter der FFGFT keine
Funktion haben \cite{Dok316}.

\section{Bezug zum Korpus}

Die verwendeten Befunde stammen aus Dok.~316 (August 2026): Weyl-Obstruktion,
Bikohärenztest, Windungszahl gegen Frequenzverhältnis, Fouriertransformation
als Projektion \cite{Dok316}. Die Überarbeitung ist unter R75 in Dok.~190
deklariert.

\section{Verifikation}

Die Aussagen dieses Dokuments sind durch drei Skripte reproduzierbar
(reine Standardbibliothek, Sollwerte als Assertions):

\begin{itemize}
  \item \texttt{s1\_periodensuche.py} -- Korrektheit der Periodenfindung,
    Skalierungsmessung (gemessener Wachstumsexponent der mittleren Ordnung
    $0{,}95$, also praktisch $O(N)$), Primfaktorzerlegung der
    $\sigma$-Nenner, Verfahrensvergleich.
  \item \texttt{s2\_grenzen.py} -- Weyl-Obstruktion gegen äquidistantes
    Register, Abgrenzung gegen Shors Aufgabe, Ressourcenabschätzung,
    Windungszahl-Probe.
  \item \texttt{s3\_thermik\_zustandsraum.py} -- thermische Besetzung nach
    Trägerfrequenz, Kühlkosten nach Carnot, Dimensionsvergleich der
    Zustandsräume, Gottesman-Knill.
\end{itemize}

\section*{Quellenverzeichnis}

\begin{thebibliography}{99}

\bibitem{Barenco1995}
Barenco et al., ``Elementary gates for quantum computation'',
\emph{Phys. Rev. A} 52, 3457 (1995).
Zeigt, dass Quantengatter elementare arithmetische Operationen
(Addition, Multiplikation) ausführen, die in ihrer Komplexität
klassischen Operationen entsprechen. Die Superposition ermöglicht
paralleles Anlegen, aber nicht paralleles Verarbeiten der Arithmetik.

\bibitem{Vedral1996}
Vedral, Barenco, Ekert, ``Quantum networks for elementary arithmetic
operations'', \emph{Phys. Rev. A} 54, 147 (1996).
Klassische Arbeit, die die Komplexität von Quanten-Addierern und
Multiplizierern analysiert und den hohen Aufwand für die modulare
Exponentiation dokumentiert.

\bibitem{Pavlidis2014}
Pavlidis, A., \& Gizopoulos, D., ``Fast quantum modular exponentiation
architecture for Shor's factoring algorithm'',
\emph{Quantum Information Processing}, Vol. 13, pp. 2899-2924 (2014).
DOI: 10.1007/s11128-014-0818-1.

\bibitem{Ragavan2018}
Ragavan, S., ``An efficient quantum modular exponentiation circuit for
Shor's algorithm'', \emph{Proceedings of SPIE} (2018).
DOI: 10.1117/12.2310178.

\bibitem{Patoary2021}
Patoary, M. N., et al., ``A DFT-based quantum modular exponentiation
algorithm'', \emph{2021 IEEE International Symposium on Circuits and
Systems (ISCAS)}.
DOI: 10.1109/ISCAS51556.2021.9401143.

\bibitem{Ekera2023}
Ekerå, M., \& Gidney, C., ``Quantum factoring in the noisy
intermediate-scale quantum era'', \emph{arXiv:2302.06647} (2023).
Zeigt, wie klassische Precomputation die Quantenschaltung entlastet und
beweist, dass Qubits die Arithmetik nicht intrinsisch beschleunigen.

\bibitem{Nielsen2010}
Nielsen, M. A., \& Chuang, I. L., \emph{Quantum Computation and Quantum
Information}, Cambridge University Press (2010), Kapitel 5.3.
Standardlehrbuch, das die modulare Exponentiation als den
kostenintensivsten Teil von Shors Algorithmus identifiziert.

\bibitem{Jozsa2003}
Jozsa, R., \& Linden, N., ``On the role of entanglement in
quantum-computational speed-up'', \emph{Proc. R. Soc. Lond. A} 459,
2011 (2003).
Zeigt Verschränkung als notwendig für exponentiellen Vorteil bei reinen
Zuständen.

\bibitem{Gidney2021}
Gidney, C., \& Ekerå, M., ``How to factor 2048 bit RSA integers in 8
hours using 20 million noisy qubits'', \emph{Quantum} 5, 433 (2021).
DOI: 10.22331/q-2021-04-15-433.
Ressourcenabschätzung für RSA-2048 auf fehlerkorrigierten
Quantencomputern.

\bibitem{Dok316}
Dok. 316 (August 2026).
Bikohärenztest an 2469 Riemann-Nullstellen, Weyl-Obstruktion,
Windungszahl gegen Frequenzverhältnis, Fouriertransformation als
Projektion.

\end{thebibliography}